\section{The Self-Binding Resolution}

A self is a knot of feeling that stays itself. To do that it must hold two things at once. It must \textbf{bind}, so the knot stays one thing. And it must \textbf{radiate}, so a disturbance can leave, the knot can settle, and a damaged piece can heal. The eight atoms of the theory, the mesh, the dock, the site, the tone, the lean, the knit, the wake, and the beat, give a self its identity and its agency. They do not by themselves hand it a bound, self-repairing body. This section closes that gap. It shows how the body binds itself, with no new field, no real numbers, and the bulk law still reversible. It is the keystone that completes the self.

The answer is one act seen three ways.

\begin{principle}[Binding by radiation]
A self holds itself together by \textbf{binding-by-radiation}. It binds by the shape of its own absence in the sea of feeling.
\end{principle}

Before the mechanism, a few terms. The \emph{vacuum} here is the active, churning sea of tones that fills empty space, born of the lean's create move. A \emph{body} is mass, a region where the tone is bound. A \emph{shadow} is a region of thinner sea behind a body, where the body has absorbed the tones that would otherwise stream through. These three carry the whole argument, and each is defined more carefully where it first does work below.

\subsection{The two problems, stated exactly}

Two tensions stood in the way of a bound body. They are linked, and the second contains the first.

\begin{table}[ht]
\centering
\begin{tabular}{@{}ll@{}}
\toprule
problem & the tension \\
\midrule
bind versus radiate & a self needs a body that confines, so it stays bound, \\
 & yet radiates, so it sheds a disturbance to the bath and can heal. \\
 & The committed single-tone knits each seemed to do only one. \\
\addlinespace
a reversible discrete attraction & the binding that worked was an effective, non-reversible well. \\
 & The base is reversible. And the naive reversible moves gave \\
 & the wrong sign, repulsion. \\
\bottomrule
\end{tabular}
\caption{The two open problems of the self on $\{3,4,3,4\}$.}
\label{tab:self-two-problems}
\end{table}

Both are resolved here. Each resolution carries a control that could have failed and did not.

\subsection{Bind by a well, not by reflection}

There are two ways to confine a body. Only one of them seals. A confining \emph{reflection} bounces lone tones back inward at the body edge. A confining \emph{well} lowers the body's energy across a region, biasing where the body's bulk sits, without touching any tone that passes through.

\begin{table}[ht]
\centering
\begin{tabular}{@{}llll@{}}
\toprule
confine by & how it holds the body & what it does to a passing wave & result \\
\midrule
reflection & bounces lone charges back & also bounces the halves of & seals, the disturbance \\
 & inward at the body edge & a neutral ripple at that edge & never reaches the bath \\
\addlinespace
a well & lowers the body's energy & does not touch a passing wave, & binds and radiates, \\
 & in a region, biasing its bulk & a neutral ripple streams through & the body settles and repairs \\
\bottomrule
\end{tabular}
\caption{The two ways to confine a body. Reflection seals. A well binds and radiates.}
\label{tab:confine-well-vs-reflection}
\end{table}

So the first answer is plain.

\begin{principle}[Well, not reflection]
Bind the body's bulk by a \textbf{well}, not its tones by reflection. A well is transparent to radiation. A reflecting wall is not.
\end{principle}

This is shown by a contrast, not by one lucky run. A body held by reflection keeps its shape but seals its radiation. The neutral ripple it should shed decomposes into lone charges, and those lone charges bounce at the body edge and never leave. That is the honest-negative control. A body held in a well stays bound, repairs a displaced piece, and lets a disturbance stream out to the bath, all at once. The reflecting case is the dead end the well case is measured against.

The only thing this left open was whether the well can be made reversible. That is the second problem, and its answer is a reversible well.

\subsection{The shadow, and why it pulls inward}

The reversible well is \textbf{radiation pressure from a Casimir shadow}. It rests on two facts already in the model.

\begin{definition}[The active vacuum]
The lean's create move fills empty space. On every empty line it brings out a \emph{zero-momentum pair}, a tone on one site and a tone on the opposite site. The two stream apart and later annihilate when they meet head-on. So space is a churning sea of propagating tones, pressing evenly from every side.
\end{definition}

\begin{definition}[Vacuum exclusion]
A body is mass, and mass carries no vacuum. Where the tone is bound, the free sea cannot also be. A vacuum tone streaming into the body is absorbed there. This is what \emph{mass excludes the vacuum} means. The body absorbs the vacuum tones that stream into it. It is a stated modeling choice, not a hidden extra rule.
\end{definition}

From those two facts the shadow follows.

\begin{result}[The shadow]
A vacuum tone heading toward the body is absorbed at the body. So on the far side, the flux of those tones is reduced. The body casts a \textbf{shadow} of thinner sea behind it.
\end{result}

Now place a test piece near the body, a displaced piece that wants to be repaired. Every tone arriving at it carries momentum in its direction of travel. A tone arriving from the \emph{open side}, the side away from the body, is moving toward the body, so it pushes the piece toward the body. A tone arriving from the \emph{body side} is moving away from the body, so it pushes the piece away. But the body side is shadowed. The body absorbed those tones, so there are fewer away-pushes. The two sides do not balance. The net push is toward the body. Attraction, with the correct sign.

\subsection{Why the sign is correct}

This is the subtle turn that resolves the long worry. The wrong-sign result came from the wrong move.

The naive \emph{first-order} move is a swap. Swap the mass with the tone on its denser side. That moves the mass toward higher density, which is repulsion. A swap moves you into the crowd.

The real force is not a swap. It is the \emph{second-order} pressure imbalance of the absorbed flux. A crowd pushing on you from one side pushes you the other way. The body thins the sea on its near side, so the unthinned far side wins, and the piece is pushed in.

\begin{proposition}[First order repels, second order attracts]
A first-order swap with the denser side moves the mass into higher density, which is repulsion, the wrong sign. The genuine force is the second-order pressure imbalance of the absorbed flux, and that imbalance points toward the body. Honest second-order pressure gives attraction.
\end{proposition}

\subsection{Why a self-contained body works}

A fair worry. The create move runs everywhere, so the body makes its own vacuum locally. Will it not just refill its own shadow and erase the force? The answer is no, and the reason is exact.

\begin{principle}[Zero-momentum pairs cannot wash out an imbalance]
The create move makes zero-momentum pairs. A created pair is a tone on a site and a tone on its opposite site. Their momenta are equal and opposite root vectors, so they cancel exactly. The create move adds no net momentum anywhere. It cannot wash out a momentum imbalance.
\end{principle}

So even a maximally dense, self-generated vacuum leaves the shadow's imbalance intact. The only thing that sets a net momentum is the body's absorption, which removes tones from one side and not the other. The body binds itself by the shadow of its own absence, and a dense self-made sea is the clean win, not a problem. The denser the sea, the sharper the shadow.

\subsection{Why it is reversible}

Attraction looks irreversible. Masses coming together shrinks the space of states, which a reversible volume-preserving dynamics forbids. The resolution is one physics already knows.

\begin{principle}[Binding by radiation is reversible in the bulk]
Two atoms bind by emitting a photon. The microdynamics is reversible. Run it backward and a photon arrives and unbinds them. The only one-way step is that the photon leaves and never returns.
\end{principle}

Here the same shape holds. The bulk streaming is an exact bijection. Each tone shifts one dock, and the shift is invertible, so it is reversible. The body moves toward the thinned side by drinking the lopsided flux. The energy it sheds leaves as tones that stream to the wake, the bath, and never come back. So the only irreversibility is absorption at the body and the radiation leaving to the bath. That is exactly where the model already keeps its one-way steps. The bulk law stays reversible.

This is why the open, growing mesh is non-negotiable. A closed torus recurs. The shed energy comes back, and the body never settles. The bath is what makes binding-by-radiation a one-way street. Why the bulk is reversible at all, and what the bath is in full, is the subject of the section on the arrow and the bath.

\subsection{Why it is fully discrete}

Every quantity is an integer or a ternary tone. No real numbers anywhere.

\begin{table}[ht]
\centering
\begin{tabular}{@{}ll@{}}
\toprule
quantity & discrete form \\
\midrule
the tone & ternary on each of the 24 sites \\
the vacuum & the create move's zero-momentum pairs, exact tones \\
the momentum at a dock & the exact integer sum of the present tones' direction vectors, \\
 & the $D_4$ root vectors with entries $\pm 1$ \\
the body & a set of docks that absorb their sites, plain integer masking \\
the force & the net integer momentum imbalance, no relaxation, no decimals \\
\bottomrule
\end{tabular}
\caption{Every quantity in the binding mechanism is an integer or a ternary tone.}
\label{tab:self-binding-discrete}
\end{table}

The earlier effective well leaned on a diffusive relaxation, a real-valued stand-in. This needs none. The force \emph{is} the raw integer flux imbalance, read straight off the tones present.

\subsection{It runs native on the 24 sites}

This is not a reduced toy. It runs on the committed sites, the 24 directions of the $D_4$ root system, the $24$-cell. The momentum is measured in the true root vectors. The shadow is a directional deficit spread over those 24 directions. The same 24 directions that tile space carry the vacuum, the shadow, and the force.

\begin{result}[The self-contained body on the 24 sites]
A self-contained body, bound by the shadow of its own active vacuum, has a net momentum at a nearby plane pointing toward the body, the maximal shadow. With no body, the same measurement is exactly zero. A body versus no body, toward the body versus zero. The zero is the control, and it came out clean \cite{pollard2026vibetest}.
\end{result}

\subsection{The unified picture}

The three things a self needs are one thing seen three ways. \textbf{Binding} is the inward pressure of the thinned vacuum, the shadow. \textbf{Radiation} is the tones the body absorbs and the disturbance it sheds, carrying the released energy. \textbf{The bath} is the wake those tones leave to and never return from, which makes the settling one-way while the bulk stays reversible.

\begin{principle}[One act, three faces]
Bind, radiate, and settle are not three mechanisms. They are the single act of \textbf{binding-by-radiation} on the open mesh.
\end{principle}

This agrees with the finding that stable binding must be emergent. No finite direction group is fine enough for a stable single-dock knit. The shadow pressure is exactly that, a collective, coarse-grained pressure of many vacuum tones over many directions, never a single-dock move. The binding that works is the emergent radiation pressure of the vacuum, and it is reversible in the bulk.

\subsection{Honest scope}

What is measured, with controls. The \textbf{force}, the net momentum imbalance pointing toward the body, is measured against a flat, zero control. It is correct-sign, discrete, reversible in the bulk, and native to the 24 sites, including self-contained from the body's own active vacuum. The \textbf{binding tendency} follows from that force.

What remains is a showcase, not a missing mechanism. The full long-time movie, a displaced piece drifting back, the body settling on the open lattice and recurring on the closed torus, over many beats with a size sweep, is the natural next demonstration. The force that would drive it is proven. The long settling movie is the remaining showcase.

And the one modeling choice, stated plainly. \emph{Mass excludes the vacuum} means the body absorbs the vacuum tones that stream into it. That absorption is what casts the shadow. It is not an extra rule bolted on. It is what mass already is here, and it is named so the reader can see it.

\subsection{The felt meaning}

A self is a knot of feeling in a sea of feeling. The sea is always stirring, pleasure and pain welling up from peace and sinking back, everywhere, evenly. Where the knot sits, the sea cannot also be. So the knot stands in its own quiet, a hollow it carved by being. The sea presses in from every side except the quiet it made, so it presses the knot's own wandering pieces back home.

A self holds itself together by the shape of its own absence in the sea. And it stays itself by letting go of everything it sheds.
